\section{Discussion}
\label{sec:discussion}

\subsection{What the evidence supports}

The split is clean and it is about labels. \textbf{Given roughly 400
labelled examples for your task, a 1.7B model with an 87\,MB adapter beats
a commercial typed-decision API} on a hardened routing benchmark, runs on
your own hardware, returns full-precision probabilities, and is
deterministic. \textbf{Without labels, it does not come close}: $0.290$ on
Banking77 against roughly $0.820$.

Both halves matter. The first says the capability is not gated behind
scale or a vendor. The second says the vendor's actual product --- working
on a schema you have never labelled --- is a real and unsolved thing.

\subsection{Why the data mattered more than the model}

The most useful finding here is negative: a strong encoder with a clever
output format scores below chance untrained. Nothing about scoring options
at marker positions confers generalization. What produced it was the
composition of the training mixture, and in particular a property we could
manufacture --- menu-size diversity --- rather than one we had to find.

This has a practical consequence. Our mixture is 234 \choice{} tasks
against 34 \noul{} and 11 \score{}, and the result pattern follows exactly:
\choice{} transfers, the other two barely do. That is a statement about
279 tasks scraped from one source, not about the architecture. The
fastest route to a better general model is a more balanced typed-decision
corpus, and assembling one is beyond what a single project can do.

\subsection{Limitations}

\paragraph{The ordinal primitive is weak.} \texttt{helpsteer} beats chance
but only draws level with predicting the most common label. The mixture
has 11 \score{} tasks, mostly three-level sentiment. We searched twice for
more; ordinal scales with five or more levels, in domains that are not
response quality, are scarce in accessible form. We deliberately excluded
LLM-response-rating corpora because they are the held-out task in all but
name, and training on a near-twin would move the number without moving the
ability.

\paragraph{No ordinal or rejection head.} \score{} is a softmax over
levels; we use the ordering only in rendering and in an expectation-based
readout. Rank-consistent objectives~\cite{corn} and adaptive rejection
boundaries~\cite{adb} are the obvious next step and are not implemented.

\paragraph{Zero-shot on a real task is much worse than on benchmarks.}
$21.9\times$ chance on the held-out suite coexists with intent $0.601$ and
a clinical-gate recall of $0.111$ on the routing task. Safety gates in
particular appear to need task supervision, and we would not deploy any
zero-shot gate on this evidence.

\paragraph{Synthetic evaluation.} The routing task is handwritten and
systematically composed, not real traffic. It is evidence a method works,
not a claim about any production distribution.

\paragraph{One-day black-box baseline.} All reference figures are
measurements of a hosted endpoint on a single day from one network
location.

\paragraph{Untuned comparisons.} The LoRA-versus-full-fine-tune result
compares one hyperparameter setting each. We believe the direction and
not the magnitude.

\subsection{Future work}

The experiment that would most change these conclusions is a balanced
typed-decision corpus with real \score{} and \noul{} diversity at the
scale Flan suggests. Beyond that: adapter composition, where a task
adapter is initialised from a general adapter, which is the
parameter-efficient analogue of the $+36$-point intermediate-task result;
rank-consistent ordinal heads; and calibrated rejection for out-of-scope
inputs.

\subsection{Conclusion}

Typed decisions do not need a large model, and they do not need
generation. They need the right output format, an attention mask that
keeps questions independent, and --- overwhelmingly --- training data
whose \emph{shape} matches deployment. Of those three, only the last was
hard, and it is the one that transfers least from the literature to
practice.
